\documentclass[11pt]{article}
\usepackage[margin=1in]{geometry}
\usepackage{amsmath,amssymb,booktabs,longtable,array,graphicx,pdflscape,float,placeins}
\usepackage{microtype,setspace,xcolor,hyperref}
\hypersetup{colorlinks=true,linkcolor=blue!55!black,urlcolor=blue!55!black}
\setlength{\parindent}{0pt}
\setlength{\parskip}{0.55em}
\onehalfspacing
\providecommand{\ResultRoot}{../results/nine_cell_design/test_run}
\providecommand{\RunLabel}{TEST RUN -- NOT FOR INFERENCE}
\providecommand{\RunAbstractNotice}{This short run validates data alignment, samplers,
aggregation, outputs, and rendering; it is not an empirical result.}
\providecommand{\PriorCaptionNotice}{Short test chains are shown only to verify the
reporting path.}
\providecommand{\ScopeNotice}{This validation run does not establish production
convergence or robustness.}
\title{\textbf{Competition and Inflation under HSA Demand}\\
\large Nine-Cell State-Space Implementation Report}
\author{}
\date{\RunLabel}

\begin{document}
\maketitle
\begin{abstract}
This report is generated by the executable implementation of the design in
\texttt{design.tex}. It estimates the pre-declared $3\times3$ price--activity
matrix using a measurement-only business-demography cut, paired quarterly annualized
and exactly aggregated year-over-year likelihoods, and a secondary full-joint
information-flow sensitivity. The displayed run is explicitly labeled above.
\RunAbstractNotice
\end{abstract}
\clearpage
\tableofcontents
\clearpage

\section{Status and interpretation boundary}
The primary coefficients are conditional varying-coefficient associations. The code
does not relabel observed firm counts as theoretical effective varieties, does not
equate aggregate inverse markup with structural marginal cost, and does not treat
full-joint feedback as an endogeneity correction. Missing benchmark-admissibility
conditions remain visible as failed or conditional gates.

\section{Data and executable design}
The common sample contains 124 quarterly endpoints from 1982Q1 through 2012Q4.
Inflation is $400\Delta\log P_t$ for PPI, headline CPI, and core CPI. The three
activity measures are the log inverse-markup proxy, the frozen BN output gap, and the
negative unemployment gap. Annual firm counts enter at Q4 only. Quarterly QCEW
establishments provide the business-demography indicator. All nine cells share one
complete-case sample; cell-specific dropping is forbidden by the loader.

The primary state posterior is the modular cut
\[
p_{cut}(S,\beta_c)=p(S\mid C)p(\beta_c\mid\pi_c,X_c,S).
\]
The full-joint E2 result is secondary. Its slow and fast state paths are sampled with
the two conditional FFBS blocks implied by the bilinear term. For YoY, the entire
quarterly conditional mean is premultiplied by $A_4$ and the likelihood uses
$G_4=A_4A_4'$.

\begin{landscape}
\section{Nine-cell primary coefficients}
\begingroup
\setlength{\tabcolsep}{3pt}
\scriptsize
\input{\ResultRoot/tables/nine_cell_coefficients.tex}
\endgroup
\end{landscape}

\section{Identification before interpretation}
\begin{center}
\input{\ResultRoot/tables/evidence.tex}
\end{center}

\begin{figure}[H]
\centering
\includegraphics[width=\textwidth]{\ResultRoot/figures/prior_posterior_cell1.png}
\caption{Cell 1 prior and primary modular-cut posterior. \PriorCaptionNotice}
\end{figure}

\section{Competition-state measurement}
\begin{figure}[H]
\centering
\includegraphics[width=\textwidth]{\ResultRoot/figures/competition_state.png}
\caption{Measurement-only slow, stationary, and total business-demography coordinate.}
\end{figure}

\IfFileExists{\ResultRoot/tables/measurement_sensitivity.tex}{
\subsection{Measurement-gate diagnosis and alternative state laws}
\begin{center}
\small
\input{\ResultRoot/tables/measurement_sensitivity.tex}
\end{center}
\begin{figure}[H]
\centering
\includegraphics[width=\textwidth]{\ResultRoot/figures/measurement_diagnostics.png}
\caption{Observed-series alignment, annual-only versus augmented uncertainty, and
the frozen alternative measurement specifications. The fast state is retained for
substantive interpretation only when every displayed specification satisfies the
pre-declared information gate.}
\end{figure}
}{}

\section{Time-varying slope}
\begin{figure}[H]
\centering
\includegraphics[width=\textwidth]{\ResultRoot/figures/slope_path_cell1.png}
\caption{Cell 1 Phillips-curve slope path under the primary cut E2 model.}
\end{figure}

\section{Cut versus full joint}
\begin{figure}[H]
\centering
\includegraphics[width=\textwidth]{\ResultRoot/figures/cut_joint.png}
\caption{Primary cut and secondary full-joint posterior summaries.}
\end{figure}
\begin{center}
\input{\ResultRoot/tables/cut_joint.tex}
\end{center}

\begin{landscape}
\section{Paired QoQ and exact YoY transformation}
\begingroup
\setlength{\tabcolsep}{2.5pt}
\scriptsize
\input{\ResultRoot/tables/qoq_yoy.tex}
\endgroup
\end{landscape}

\IfFileExists{\ResultRoot/tables/focus_robustness.tex}{
\section{Persistent-error, low-frequency, and slow-only follow-up}
The persistent-error specification uses a stationary AR(1) quarterly inflation
disturbance. The low-frequency specification adds the pre-declared latent AR(1)
inflation nuisance symmetrically across E0--E2. When the measurement gate fails, the
slow-only specification is the reporting fallback rather than an ex post replacement
for the visible E2 result.

\begin{figure}[H]
\centering
\includegraphics[width=\textwidth]{\ResultRoot/figures/focus_robustness.png}
\caption{Focused $\kappa_1$ posterior summaries for Cells 2, 5, 6, and 9.}
\end{figure}

\begin{center}
\small
\input{\ResultRoot/tables/focus_robustness.tex}
\end{center}

\begin{landscape}
\subsection{Global robustness-shift rule}
\begin{center}
\scriptsize
\input{\ResultRoot/tables/inflation_robustness.tex}
\end{center}
\end{landscape}
}{}

\section{Economic magnitudes}
\begin{center}
\small
\input{\ResultRoot/tables/effects.tex}
\end{center}
The table reports an IQR-by-IQR slope-change contribution in annualized inflation
percentage points. It enables comparison in common outcome units without treating raw
coefficients on differently scaled activity variables as comparable structural objects.

\section{Auxiliary HSA benchmark}
\begin{center}
\input{\ResultRoot/tables/benchmark.tex}
\end{center}
\begin{figure}[H]
\centering
\includegraphics[width=\textwidth]{\ResultRoot/figures/locality.png}
\caption{Frozen local-support rule based only on the measurement cut.}
\end{figure}
The benchmark is conditional on the unit-scale $b_x=1$ calibration, maintained
activity orthogonality, a GDP-deflator expectation proxy for PPI, and a calibrated
reference markup. Compatibility is not a probability that HSA is true.

\section{What this run does not establish}
This run does not establish causal competition effects, a structural activity equation,
real-time predictive validity, or a fully admissible PPI benchmark. \ScopeNotice

\end{document}
